\documentclass[a4paper,12pt,titlepage,finall]{article}

\usepackage[T1,T2A]{fontenc}     % форматы шрифтов
\usepackage[utf8x]{inputenc}     % кодировка символов, используемая в данном файле
\usepackage[russian]{babel}      % пакет русификации
\usepackage{tikz}                % для создания иллюстраций
\usepackage{pgfplots}            % для вывода графиков функций
\usepackage{geometry}		 % для настройки размера полей
\usepackage{indentfirst}         % для отступа в первом абзаце секции
\usepackage{multirow}            % для таблицы с результатами

% выбираем размер листа А4, все поля ставим по 3см
\geometry{a4paper,left=30mm,top=30mm,bottom=30mm,right=30mm}

\setcounter{secnumdepth}{0}      % отключаем нумерацию секций

\usepgfplotslibrary{fillbetween} % для изображения областей на графиках

\begin{document}
% Титульный лист
\begin{titlepage}
    \begin{center}
	{\small \sc Московский государственный университет \\имени М.~В.~Ломоносова\\
	Факультет вычислительной математики и кибернетики\\}
	\vfill
	{\Large \sc Отчет по заданию №1}\\
	~\\
	{\large \bf <<Методы сортировки>>}\\ 
	~\\
	{\large \bf Вариант 4}
    \end{center}
    \begin{flushright}
	\vfill {Выполнил:\\
	студент 104 группы\\
	Заворочаев~Л.~В.\\
	~\\
	Преподаватель:\\
	Гуляев~Д.~А.}
    \end{flushright}
    \begin{center}
	\vfill
	{\small Москва\\2026}
    \end{center}
\end{titlepage}

% Автоматически генерируем оглавление на отдельной странице
\tableofcontents
\newpage

\section{Постановка задачи}

В рамках поставленной задачи требовалось сравнить два метода сортировки по числу сравнений и перестановок элементов. В данном варианте исследовались сортировка Хоара и сортировка вставками. Сортировки проводились в порядке неубывания элементов, представляющих собой числа двойной точности (double).

Результаты исследования были получены с помощью тестов, в ходе которых каждой функции сортировки на вход подавались массивы размера 10, 100, 1000, 10000. Для каждого размера подавалось 4 массива со следующими свойствами:
\begin{enumerate}
\item Массив, упорядоченный по возрастанию;
\item Массив, упорядоченный по убыванию;
\item Массив со случайной расстановкой элементов;
\item Массив со случайной расстановкой элементов. 
\end{enumerate}

Тесты проводились для функий сортировок на основе одних и тех же данных.

\newpage

\section{Результаты экспериментов}

\begin{table}[h]
\centering
\begin{tabular}{|c|c|c|c|c|c|c|c|}
    \hline
    \multirow{2}{*}{\textbf{n}} & \multirow{2}{*}{\textbf{Параметр}} & \multicolumn{4}{|c|}{\textbf{Номер сгенерированного массива}} & \textbf{Среднее} \\
    \cline{3-6}
    & & \parbox{1.5cm}{\centering 1} & \parbox{1.5cm}{\centering 2} & \parbox{1.5cm}{\centering 3} & \parbox{1.5cm}{\centering 4} & \textbf{значение} \\
    \hline
    \multirow{2}{*}{10} & Сравнения & 43 & 42 & 39 & 58 & 45.50 \\
    \cline{2-7}
                        & Перемещения & 0 & 5 & 7 & 4 & 4.00 \\
    \hline
    \multirow{2}{*}{100} & Сравнения & 771 & 770 & 834 & 848 & 805.75 \\
    \cline{2-7}
                         & Перемещения & 0 & 50 & 135 & 129 & 78.50 \\
    \hline
    \multirow{2}{*}{1000} & Сравнения & 10975 & 10974 & 12982 & 15439 & 12592.50 \\
    \cline{2-7}
                          & Перемещения & 0 & 500 & 2069 & 1992 & 1140.25 \\
    \hline
    \multirow{2}{*}{10000} & Сравнения & 143615 & 143614 & 180024 & 177805 & 161264.50 \\
    \cline{2-7}
                           & Перемещения & 0 & 5000 & 28140 & 28367 & 15376.75 \\
    \hline
\end{tabular}
\caption{Результаты работы сортировки Хоара (Quick Sort)}
\end{table}

\begin{table}[h]
\centering
\begin{tabular}{|c|c|c|c|c|c|c|c|}
    \hline
    \multirow{2}{*}{\textbf{n}} & \multirow{2}{*}{\textbf{Параметр}} & \multicolumn{4}{|c|}{\textbf{Номер сгенерированного массива}} & \textbf{Среднее} \\
    \cline{3-6}
    & & \parbox{1.5cm}{\centering 1} & \parbox{1.5cm}{\centering 2} & \parbox{1.5cm}{\centering 3} & \parbox{1.5cm}{\centering 4} & \textbf{значение} \\
    \hline
    \multirow{2}{*}{10} & Сравнения & 9 & 45 & 34 & 26 & 28.50 \\
    \cline{2-7}
                        & Перемещения & 0 & 45 & 27 & 20 & 23.00 \\
    \hline
    \multirow{2}{*}{100} & Сравнения & 99 & 4950 & 2453 & 2383 & 2471.25 \\
    \cline{2-7}
                         & Перемещения & 0 & 4950 & 2359 & 2290 & 2339.75 \\
    \hline
    \multirow{2}{*}{1000} & Сравнения & 999 & 499500 & 247146 & 258834 & 251619.75 \\
    \cline{2-7}
                          & Перемещения & 0 & 499500 & 246155 & 257843 & 250874.50 \\
    \hline
    \multirow{2}{*}{10000} & Сравнения & 9999 & 49995000 & 25030722 & 24853354 & 24972268.25 \\
    \cline{2-7}
                           & Перемещения & 0 & 49995000 & 25020731 & 24843364 & 24964773.75 \\
    \hline
\end{tabular}
\caption{Результаты работы сортировки вставками (Insertion sort)}
\end{table}

\newpage

Из приведённой таблицы видно, что сортировка на малых объёмах данных сортировка втавками выполняет меньшее число сравнений, чем быстрая сортировка, однако с увеличением размера массива, сортировка Хоара становится эффективнее сортировки вставками как по числу сравнений, так и по количеству перестановок.
\section*{Сложность алгоритмов}
\subsection*{По числу сравнений}
Сортировка Хоара:
\begin{itemize}
	\item Лучший случай: $O(n \log n)$
	\item В среднем: $O(n \log n)$
	\item Худший случай: $O(n^2)$, теоретически возможен, однако не наблюдается в тестах
\end{itemize}
Сортировка вставками:
\begin{itemize}
	\item Лучший случай: $O(n)$
	\item В среднем: $O(n^2)$
	\item Худший случай: $O(n^2)$
\end{itemize}

\subsection*{По числу перестановок}
Сортировка Хоара:
\begin{itemize}
	\item Лучший случай: $O(1)$
	\item В среднем: $O(n \log n)$
	\item Худший случай: $O(n^2)$, теоретически возможен, однако не наблюдается в тестах
\end{itemize}
Сортировка вставками:
\begin{itemize}
	\item Лучший случай: $O(1)$
	\item В среднем: $O(n^2)$
	\item Худший случай: $O(n^2)$
\end{itemize}

\newpage

\section{Структура программы и спецификация функций}

Программа написана на Си и поделена на модули main.c (Проведение тестов и вывод программы) generation.c (Генерация массивов) и sortings.c (Сортировки).

\subsection{Глобальные переменные}

\texttt{compares} — счётчик числа сравнений последний проведённой сортировки

\texttt{swaps} — счётчик числа перестановок последний проведённой сортировки

\subsection{Функции сравнения}

\texttt{int ascending(const void *a, const void *b)} — компаратор для мортировки по возрастанию

\texttt{int descending(const void *a, const void *b)} — компаратор для сортировки по убыванию

\subsection{Вспомогательные функции}

\texttt{void zeroize(void)} — обнуляет глобальные счётчики

\texttt{unsigned long long llrand(void)} — генерирует большое случайное число. Используется для генерации дробных числе в большем диапазоне

\texttt{double rand\_double(void)} — генерирует случайное число с плавающей точкой

\texttt{void copy\_array(double to\_here[], double from\_here[], int size)} — копирует один массив в другой. Используется для независимой сортировки одних и тез же данных разными способами

\texttt{void set\_tests(int tests[4][2][2], int size)} — запоминает в массиве \texttt{tests} результаты проведённых тестов

\subsection{Функции генерации массивов}

\texttt{void quick\_sort(double arr[], int begin, int end)} — функция сортировки Хоара

\texttt{int partition(double arr[], int left, int right)} — вспомогательная функция для сортировки Хоара

\texttt{void insertion\_sort(double arr[], int size)} — функция сортировки вставками

\section{Отладка программы, тестирование функций}
В ходе отладки тестировались методы сортировок, проверялось корректность сортировок путём вывода массивов небольших размеров, прошедших через функции.

\section{Анализ допущенных ошибок}
Попытка рандомизации дробных чисел путём побитовой рандомизации приводила к генерации чрезмерно больших или чрезмерно маленьких чисел, которых сложно было читать во время отладки.

\begin{raggedright}
\addcontentsline{toc}{section}{Список цитируемой литературы}
\begin{thebibliography}{99}
\bibitem{cs} Кормен Т., Лейзерсон Ч., Ривест Р, Штайн К. Алгоритмы: построение и анализ.
    Второе издание.~--- М.:<<Вильямс>>, 2005.
\end{thebibliography}
\end{raggedright}

\end{document}